\section{CKKS on Zoo}

\label{sec:ckks}
%----------------------------------------------------------------------

\subsection{Approximate Arithmetic}

\CKKS{}~\cite{cheon2017ckks} encodes a vector $\mathbf{z} \in \mathbb{C}^{N/2}$ into a polynomial ring $R_q = \mathbb{Z}_q[X]/(X^N+1)$ via the canonical embedding inverse, scaled by $\Delta$:

\begin{equation}
  m(X) = \lfloor \Delta \cdot \sigma^{-1}(\mathbf{z}) \rceil \in R_q
\end{equation}

\noindent where $\Delta = 2^{40}$ is the scaling factor and $\sigma^{-1}$ is the inverse canonical embedding. Encryption produces $\mathsf{ct} = (\mathsf{ct}_0, \mathsf{ct}_1) \in R_q^2$ under \RLWE{}.

\subsection{Homomorphic Operations}

Addition of two \CKKS{} ciphertexts at the same level is component-wise. Multiplication produces a degree-2 ciphertext that must be relinearized:

\begin{align}
  \mathsf{ct}_{\mathrm{mult}} &= (\mathsf{ct}_0^{(a)} \cdot \mathsf{ct}_0^{(b)},\; \mathsf{ct}_0^{(a)} \cdot \mathsf{ct}_1^{(b)} + \mathsf{ct}_1^{(a)} \cdot \mathsf{ct}_0^{(b)},\; \mathsf{ct}_1^{(a)} \cdot \mathsf{ct}_1^{(b)}) \\
  \mathsf{ct}_{\mathrm{relin}} &= \mathsf{Relin}(\mathsf{ct}_{\mathrm{mult}}, \mathsf{evk})
\end{align}

\subsection{Rescaling and Level Management}

After multiplication, the scaling factor doubles to $\Delta^2$. Rescaling divides by $\Delta$ and drops one prime from the ciphertext modulus chain:

\begin{equation}
  \mathsf{Rescale}(\mathsf{ct}, q_\ell) \to \mathsf{ct}' \text{ at level } \ell - 1, \quad q_{\ell-1} = q_\ell / p_\ell
\end{equation}

\Zoo{} provisions a modulus chain of $L = 15$ levels by default, supporting up to 14 sequential multiplications---sufficient for a 4-layer neural network with ReLU approximations. The VM tracks levels alongside gas, rejecting operations that would exhaust the chain.

\subsection{ML Inference on Encrypted Data}

A neural-network forward pass on encrypted inputs uses \CKKS{} slot packing. For a weight matrix $W \in \mathbb{R}^{m \times n}$ and encrypted input $\mathsf{Enc}(\mathbf{x})$:

\begin{equation}
  \mathsf{Enc}(W\mathbf{x} + \mathbf{b}) = W \otimes \mathsf{Enc}(\mathbf{x}) \oplus \mathsf{Enc}(\mathbf{b})
\end{equation}

\noindent where $\otimes$ denotes plaintext--ciphertext multiplication and $\oplus$ denotes ciphertext addition. Activation functions are approximated by low-degree polynomials (degree-3 Chebyshev for ReLU) consuming one multiplicative level each.

%----------------------------------------------------------------------
